% =====================================================================
% experiment_design_phil_sci.tex
%
% Experiment-design document for the multi-agent "change of mind" /
% paradigm-shift active-inference model. Written from the joint
% perspective of a philosopher of science (Mayo / Lakatos / Kuhn
% lineage) and a computational social scientist (Centola lineage),
% with the active-inference formalism kept honest (Friston).
%
% The model has TWO interchangeable implementations in this repo:
%   (S) the structural Gaussian information-form Bayes-net   src/structural/
%   (P) the categorical epistemic POMDP                      src/pomdp/
% Every experiment below is specified against the real knobs of both.
%
% Compiles standalone:  pdflatex experiment_design_phil_sci
% (citations use an embedded thebibliography so no bibtex pass is needed;
%  keys that also live in ../paper/refs.bib are noted in comments.)
% =====================================================================
\documentclass[11pt]{article}

\usepackage[a4paper,margin=1in]{geometry}
\usepackage[T1]{fontenc}
\usepackage[utf8]{inputenc}
\usepackage{amsmath,amssymb,amsthm,mathtools,bm}
\usepackage{booktabs}
\usepackage{array}
\usepackage{microtype}
\usepackage{xcolor}
\usepackage{enumitem}
\usepackage{hyperref}
\hypersetup{colorlinks=true,linkcolor=blue!45!black,
            citecolor=blue!45!black,urlcolor=blue!45!black}

\newcommand{\code}[1]{\texttt{#1}}
\newcommand{\param}[1]{\texttt{#1}}
\newcommand{\Pib}{\boldsymbol{\Pi}}
\newcommand{\hb}{\mathbf{h}}
\newcommand{\thetastar}{\theta^{\!*}}
\newcommand{\mt}{m_t}
\newcommand{\Hzero}{\mathcal{H}_0}
\newcommand{\Hone}{\mathcal{H}_1}

\newtheorem{hyp}{Hypothesis}
\newtheorem{prediction}{Prediction}

% severe-test callout box
\newcommand{\severe}[1]{%
  \par\smallskip\noindent
  \fcolorbox{red!40!black}{red!4}{%
    \parbox{\dimexpr\linewidth-2\fboxsep-2\fboxrule}{%
      \textbf{\textcolor{red!50!black}{Severe test / falsifier.}}\ #1}}%
  \par\smallskip}

\title{\textbf{Three Severe Tests of a Multi-Agent Model of\\ Paradigm Change}\\[4pt]
\large An experiment-design document at the intersection of\\
philosophy of science and computational social science}
\author{Forest-walk synthesis \\ \small (Mayo $\cdot$ Lakatos $\cdot$ Centola $\cdot$ Friston)}
\date{2026-06-01}

\begin{document}
\maketitle

\begin{abstract}
\noindent
We design three experiments on a multi-agent active-inference model of scientific
paradigm change. The model exists in two interchangeable implementations: a
\emph{structural} Gaussian information-form Bayes-net in which a paradigm \emph{is}
a precision matrix (\code{src/structural/}), and a \emph{categorical epistemic
POMDP} in which a fixed true paradigm $\thetastar$ is read through theory-laden
likelihoods and probed by experiments chosen to minimise expected free energy
(\code{src/pomdp/}). The model's central empirical claim to date is that its
dynamics are \emph{monostable}---a single attractor at the truth---and that
apparent ``lock-in'' is \emph{self-censorship} (a cost-driven refusal to run
discriminating experiments) and \emph{paradigm-governed precision silencing},
not genuine bistability. Following Mayo's principle that a claim is only
established when it has survived a test it would probably have failed if false,
each experiment is built around an explicit falsifier. Experiment~1 probes the
\emph{ontology} (is there a second basin?) via rate-resolved hysteresis.
Experiment~2 probes the \emph{individual epistemics} (when does a research
programme seal itself?) by mapping the progressive/degenerating frontier in the
cost $\times$ conviction plane. Experiment~3 probes the \emph{collective
epistemics} (when can a structured society break the seal?) via committed-minority
tipping under varying trust topology. All manipulated variables are named code
parameters with ranges; all outcomes are existing order parameters.
\end{abstract}

\tableofcontents

% =====================================================================
\section{The model under test}
\label{sec:model}

\subsection{Two implementations of one object}

\paragraph{(S) Structural Bayes-net.}
Each agent holds a Gaussian belief over $d$ interdependent commitments in
information form $(\Pib,\hb)$ with $\Pib=\Sigma^{-1}$ and $\hb=\Pib\bm\mu$.
A \emph{paradigm} is the off-diagonal structure of $\Pib$: a dense hub block
(the Lakatosian \emph{hard core}) coupled to peripheral commitments (the
\emph{protective belt}). The canonical scenario is the $d{=}10$ phlogiston
network (\code{src/structural/phlogiston.py}) with three
\code{DISAGREEMENT\_NODES} on which phlogiston and oxygen priors conflict.
Per round the population does
\begin{equation}
\underbrace{\Pib_f=\textstyle\sum_j W_{ij}\Pib_j,\quad \hb_f=W\hb}_{\text{FUSE (trust-weighted pooling)}}
\;\longrightarrow\;
\underbrace{\Pib=\Pib_f+\textstyle\sum_k \rho_k\,\mathbf{H}_k^\top\mathbf{H}_k/\sigma_o^2}_{\text{OBSERVE (attention-gated Fisher deposit)}},
\end{equation}
where the per-channel \emph{evidential precision} $\rho_k$ is the model's
attention. Bayesian Model Reduction (BMR) prices the removal of a node by the
closed-form log Bayes factor $\Delta F$ (Savage--Dickey), and the
\emph{carry-over} $\Pib_{ab}\Pib_{bb}^{-1}\Pib_{ba}$ (Schur complement) is the
literal cost of revising the core. The order parameter is
\begin{equation}
\mt = \frac{\overline{\mu}_{\text{disagree}}(t)-\mu_{\text{phlog}}}{\mu_{\text{oxy}}-\mu_{\text{phlog}}}
\in[0,1],\qquad \mt{=}0\ \text{phlogiston held},\ \ \mt{=}1\ \text{oxygen held}.
\end{equation}

\paragraph{(P) Categorical epistemic POMDP.}
A fixed true paradigm $\thetastar\in\{0,\dots,K-1\}$ generates outcomes through a
theory-laden likelihood $A_{\text{world}}[o,\theta,a]$ (columns differ across
$\theta$: the same outcome is read differently under different theories).
Each agent infers $q(\theta)$ by exact categorical Bayes from a chosen
experiment's outcome and a trust-routed social report
$o_{\text{social}}=T\,q$ with channel fidelity \param{q\_reliability}. Experiments
$a$ are chosen by softmax over the negative expected free energy
\begin{equation}
\label{eq:efe}
-\mathrm{EFE}(a)\;=\;\underbrace{\mathbb{E}_{o}\,\mathrm{KL}\!\big[q(\theta\,|\,o,a)\,\|\,q(\theta)\big]}_{\text{epistemic (salience)}}
\;+\;\underbrace{\mathbb{E}_{o}\,C[o]}_{\text{pragmatic}}
\;+\;\beta_U\,\mathrm{bu}(a)\;-\;\underbrace{\param{cost\_scale}\cdot c[a]}_{\text{experiment effort}},
\end{equation}
with policy $q(a)=\mathrm{softmax}(\param{gamma\_policy}\cdot(-\mathrm{EFE}))$.
Order parameters returned by \code{run\_fast}: \param{mean\_qB} (mean belief in
the truth), \param{occ\_B} (fraction with MAP${=}\thetastar$), \param{p\_discrim}
(policy mass on the most-discriminating experiment), \param{chosen\_d} (mean
discriminability of chosen experiments).

\subsection{The claim being probed, and the methodological stance}

\noindent The standing empirical claim (notebooks 18--20) is:
\begin{quote}\itshape
The model is monostable---every trajectory relaxes to $\mu=\thetastar$. ``Lock-in''
is not a second attractor; it is (i) cost-driven \emph{self-censorship}, the
EFE-rational refusal of a confident agent to run an expensive, low-information
refuting experiment, and (ii) \emph{paradigm-governed precision silencing} via
the \param{core\_governance} knob. Trust-weighted social disagreement can break
the self-seal.
\end{quote}

We adopt Mayo's \emph{severe testing} criterion \cite{mayo2018severe}: a claim is
corroborated only by a procedure that \emph{probably would have produced a result
discordant with the claim, were the claim false}. Each experiment therefore
foregrounds the specific result that, if observed, refutes the standing claim.
The three experiments are ordered by level of description---ontology, individual,
collective---mirroring the Lakatosian distinction between a programme's hard core,
its problemshifts, and its competition with rivals \cite{lakatos1970methodology}.

% =====================================================================
\section{Experiment 1 --- Rate-resolved hysteresis: is the object bistable?}
\label{sec:exp1}

\subsection{Motivation}
Kuhn \cite{kuhn1962Structure} describes paradigm change as path-dependent and
gestalt-like: once a community converts, returning the evidence to its earlier
state does not return the community. That is the phenomenology of \emph{bistability}
(two attractors with a separating basin boundary). The model claims the opposite:
\emph{monostability with tracking lag}. The two are observationally identical under
a one-directional drive, which is exactly why a severe test is needed. The
discriminating instrument is a \emph{reversible} drive applied at several speeds:
a tracking lag yields a hysteresis loop whose area vanishes as the drive slows,
whereas a genuine fold leaves a residual loop in the quasi-static limit.

\subsection{Design}
We drive the population with a triangular evidence schedule. In (S) we ramp the
world commitment $\phi_{\text{true}}$ on the disagreement nodes from the
phlogiston pole to the oxygen pole and back, replacing the single \param{t\_shift}
flip by a controlled ramp of rate $r$ (pole-to-pole in $1/r$ steps). In (P) we
ramp the prior/evidence balance equivalently by scheduling which paradigm the
environment samples from. For each ramp we record the loop in $(\,\text{evidence},
\mt\,)$ space and compute the loop area
\begin{equation}
A_{\text{hys}}(r,g) \;=\; \oint \mt \, \mathrm{d}(\text{evidence}),
\end{equation}
as a function of ramp rate $r$ and the lock-in knob (\param{core\_governance}
$g$ in (S); \param{cost\_scale} and \param{gamma\_policy} in (P)). We
additionally run a \emph{basin probe}: initialise at the oxygen pole, apply a
transient phlogiston evidence pulse of amplitude $\Delta$ and duration $\tau$,
and record whether the population returns to oxygen or is captured.

\subsection{Manipulated variables}
\begin{center}
\begin{tabular}{@{}lll@{}}
\toprule
Parameter & Implementation & Levels \\
\midrule
ramp rate $r$ & both & $\{1/10,\,1/30,\,1/100,\,1/300\}$ (steps$^{-1}$); quasi-static $=1/300$ \\
\param{core\_governance} $g$ & (S) & $\{0,\ 1,\ 10,\ 100,\ 1000\}$ \\
\param{cost\_scale} & (P) & $\{0,\ 0.25,\ 0.5,\ 1.0\}$ \\
\param{gamma\_policy} & (P) & $\{1,\ 5,\ \infty\}$ ($\infty\equiv$ \param{greedy=True}) \\
pulse amplitude $\Delta$ & both & $5$ levels spanning $[0, 2(\mu_{\text{oxy}}-\mu_{\text{phlog}})]$ \\
pulse duration $\tau$ & both & $\{2,5,10,20\}$ steps \\
\bottomrule
\end{tabular}
\end{center}

\subsection{Fixed parameters}
$d{=}10$, $N{=}60$, \param{sigma\_o}${=}1.0$, \param{base\_prec}${=}1.0$,
\param{hub\_coupling}${=}0.8$, \param{mass\_prec}${=}2.5$ (S);
$K{=}2$, $\thetastar{=}1$, \param{x\_grid}${=}(0.1,0.5,1.0,2.0,3.0)$,
\param{q\_reliability}${=}0.80$ (P). $n_{\text{seeds}}{=}40$. Network:
Watts--Strogatz, \param{mean\_degree}${=}4$.

\subsection{Outcomes and analysis}
Primary outcome: the slope $\partial A_{\text{hys}}/\partial r$ near $r\to 0$ and
the intercept $A_{\text{hys}}(0,g)\equiv\lim_{r\to0}A_{\text{hys}}$ obtained by
linear (and saturating) extrapolation across the four ramp rates. Secondary:
fraction of seeds captured at horizon in the basin probe, as a function of
$(\Delta,\tau)$, giving an empirical basin boundary; escape-time distribution in
(P) under finite vs.\ infinite \param{gamma\_policy}.

\begin{hyp}[Monostable, standing claim]\label{h:mono}
$A_{\text{hys}}(0,g)=0$ for all finite $g$ and all finite \param{gamma\_policy}:
the loop area extrapolates to zero in the quasi-static limit; the basin probe
always returns to oxygen for finite $\tau$.
\end{hyp}
\begin{hyp}[Bistable, Kuhnian alternative]\label{h:bi}
$\exists\,g^{*}$ such that $A_{\text{hys}}(0,g)>0$ for $g>g^{*}$: a residual,
rate-independent loop survives; the basin probe exhibits a sharp
$(\Delta,\tau)$ capture boundary.
\end{hyp}

\begin{prediction}
From notebooks 19--20 we expect Hypothesis~\ref{h:mono} to hold for all
\emph{finite} settings, with $A_{\text{hys}}\propto r$ (pure tracking lag), and a
\emph{degenerate} fold appearing only in the singular limit
\param{gamma\_policy}$=\infty$ (\param{greedy=True}), where escape probability is
exactly zero. \param{core\_governance} should rescale the relaxation rate (hence
the finite-$r$ loop area) without creating a residual loop.
\end{prediction}

\severe{If the loop area extrapolated to $r\to 0$ is statistically
indistinguishable from zero (one-sided test, $\alpha{=}0.05$, against the
$n_{\text{seeds}}{=}40$ spread) at every finite $g$, monostability survives a
test it would have failed had a second basin existed. Conversely, a residual
loop at any finite $g$ \emph{refutes} the standing monostability claim. Reporting
$A_{\text{hys}}$ at a single ramp rate is \emph{not} a severe test and is
explicitly disallowed.}

% =====================================================================
\section{Experiment 2 --- The degenerating-programme frontier:
when does the population stop running the crucial experiment?}
\label{sec:exp2}

\subsection{Motivation}
Lakatos \cite{lakatos1970methodology} distinguishes \emph{progressive} research
programmes (which keep risking and corroborating novel predictions) from
\emph{degenerating} ones (which expend their effort on post-hoc self-protection).
In this model that distinction is not about belief but about \emph{behaviour}: a
degenerating programme is one whose agents stop choosing the discriminating
experiment. By the Duhem--Quine logic, the ``crucial experiment'' is never
logically forced; here it is \emph{economically} avoided---for a confident agent
the refuting experiment has low expected information gain (Eq.~\ref{eq:efe},
epistemic term) and nonzero cost, so EFE rationally deselects it. The key
non-obvious claim (notebook 19) is that this self-censorship appears \emph{even at
zero cost} above a conviction threshold. Experiment~2 maps the frontier and gives
the zero-cost null its best chance to survive.

\subsection{Design}
Fully crossed factorial in (P) over experiment cost, initial conviction in the
\emph{wrong} paradigm, and policy precision; replicated as a robustness check in
(S) via \param{experiment\_bias} and \param{core\_governance}. We track the
\emph{behavioural} order parameters \param{p\_discrim} and \param{chosen\_d} over
the horizon, alongside the \emph{doxastic} \param{mean\_qB}. The diagnostic of a
degenerating programme is the joint event
\[
\big(\,\param{chosen\_d}(t)\ \text{declining}\,\big)\ \wedge\ \big(\,\param{mean\_qB}(t)<\tfrac12\ \text{stuck}\,\big),
\]
i.e.\ the population both \emph{stops looking} and \emph{stays wrong}.

\subsection{Manipulated variables}
\begin{center}
\begin{tabular}{@{}lll@{}}
\toprule
Parameter & Implementation & Levels \\
\midrule
\param{cost\_scale} & (P) & $\{0,\ 0.1,\ 0.25,\ 0.5,\ 1.0,\ 2.0\}$ \\
initial conviction $D_0$ on wrong $\theta$ & (P) & $\{0.5, 0.6, 0.7, 0.8, 0.9, 0.95\}$ \\
\param{gamma\_policy} & (P) & $\{1,\ 5,\ \infty\}$ \\
\param{cost\_kind} & (P) & $\{$\code{discriminability}, \code{power}, \code{vector}$\}$ \\
\param{experiment\_bias} & (S) & $\{0,\ 0.25,\ 0.5,\ 0.75,\ 1.0\}$ \\
\param{core\_governance} $g$ & (S) & $\{0,\ 1,\ 10,\ 100\}$ \\
\bottomrule
\end{tabular}
\end{center}

\subsection{Fixed parameters}
$K{=}2$, $\thetastar{=}1$, \param{theta\_vals}${=}(0,1)$, \param{beta\_U}${=}0$
(switched off here so the effect is purely epistemic$+$cost; \param{beta\_U} is
varied separately in a supplementary sweep), \param{n\_o}${=}5$,
\param{x\_grid}${=}(0.1,0.5,1.0,2.0,3.0)$, \param{q\_reliability}${=}0.80$,
$N{=}60$, horizon ${=}200$, $n_{\text{seeds}}{=}40$.

\subsection{Outcomes and analysis}
Per cell we summarise (i) the terminal \param{chosen\_d}, (ii) its time-slope over
the second half of the run, and (iii) terminal \param{mean\_qB}. The
\emph{self-censorship boundary} is the contour in the $(\param{cost\_scale},D_0)$
plane where terminal \param{chosen\_d} drops below a fixed fraction (say $0.5$) of
its cost-zero, low-conviction value. This contour \emph{is} the
progressive/degenerating frontier, drawn quantitatively.

\begin{hyp}[Intrinsic self-censorship]\label{h:intrinsic}
At \param{cost\_scale}${=}0$ there exists a conviction threshold
$D_0^{c}\approx 0.7$ above which \param{chosen\_d} collapses under
\param{gamma\_policy}$=\infty$: self-censorship is not purely a cost artifact.
\end{hyp}
\begin{hyp}[Cost-only self-censorship, null]\label{h:costonly}
At \param{cost\_scale}${=}0$, \param{chosen\_d} never collapses for any $D_0$ or
\param{gamma\_policy}: self-censorship requires nonzero experiment cost.
\end{hyp}

\begin{prediction}
Hypothesis~\ref{h:intrinsic} holds under greedy policies; under sampled
(\param{gamma\_policy} finite) policies the population almost always eventually
runs the crucial experiment and escapes, so the collapse is conviction- and
$\gamma$-gated, not cost-gated. The frontier should shift toward lower $D_0$ as
\param{cost\_scale} rises. The (S) replication should show \param{rho\_k} on the
disagreement channels collapsing along a homologous contour in
$(\param{experiment\_bias}, g)$.
\end{prediction}

\severe{Hypothesis~\ref{h:costonly} is the null we are obliged to give a fair
hearing: run the complete $D_0\times\param{gamma\_policy}$ sweep \emph{at exactly
\param{cost\_scale}${=}0$}. If \param{chosen\_d} stays high everywhere, the
``epistemic self-censorship'' story is refuted and the mechanism is merely
pragmatic cost. Multiple realisability is itself a test: a self-censorship
contour that appears in (P) but \emph{not} in the (S) \param{rho\_k} silencing
(or vice versa) indicates an implementation artifact rather than a model
mechanism, and must be reported as a failure to corroborate.}

% =====================================================================
\section{Experiment 3 --- Committed-minority tipping:
when can a structured society break the self-seal?}
\label{sec:exp3}

\subsection{Motivation}
The first two experiments treat the population as effectively homogeneous. The
computational-social-science question is whether the \emph{topology} of trust
decides the outcome. Centola \cite{centola2018how,centola2010spread} distinguishes
\emph{simple} contagion (one exposure suffices; conversion roughly linear in the
seed fraction) from \emph{complex} contagion (conversion needs reinforcement from
several trusted sources, producing a sharp critical fraction $f^{*}$ that depends
on network structure---famously near $25\%$, and lowered by wide bridges). Here a
committed minority is a \emph{vanguard} that keeps running the crucial experiment
and emits the correct report; the legitimate rescue channel is the structural EFE
\emph{social} drive \param{social\_weight} $\beta_s$, which re-opens a self-sealed
channel when a \emph{trusted peer who has actually looked} disagrees
(\code{notes/efe\_attention.tex}). Notebook~10 already hints that an
\emph{isolated} high-centrality vanguard overshoots and is discounted---a
prediction we now test directly.

\subsection{Design}
Take a mainstream population pre-sealed by the Experiment-2 settings
(\param{cost\_scale}${>}0$ or $g{>}0$ above the self-censorship frontier). Inject
a vanguard fraction $f$ of agents pinned to the truth who keep running the
discriminating experiment, and vary (a) $f$, (b) the trust topology, (c) the
vanguard's \emph{placement}, and (d) the rescue strength. We measure whether and
how fast the mainstream flips.

\subsection{Manipulated variables}
\begin{center}
\begin{tabular}{@{}lll@{}}
\toprule
Parameter & Implementation & Levels \\
\midrule
vanguard fraction $f$ & both & $\{0,0.05,0.10,0.15,0.20,0.25,0.30,0.40,0.50\}$ \\
trust topology \param{kind} & both & $\{$Watts--Strogatz, ring lattice, random (ER), scale-free (BA)$\}$ \\
vanguard placement & both & $\{$random, high-centrality, clustered-behind-wide-bridge$\}$ \\
\param{mean\_degree} & both & $\{2,4,8\}$ \\
\param{social\_weight} $\beta_s$ & (S) & $\{0,\ 0.5,\ 1.0,\ 2.0\}$ \\
\param{q\_reliability} & (P) & $\{0.6,\ 0.8,\ 0.95\}$ \\
\bottomrule
\end{tabular}
\end{center}

\subsection{Fixed parameters}
Mainstream sealed at \param{cost\_scale}${=}1.0$ / $g{=}100$ (above the
Experiment-2 frontier); $N{=}120$ (larger, to resolve $f$ in $5\%$ steps);
$\thetastar{=}1$; horizon ${=}300$; $n_{\text{seeds}}{=}50$;
\param{gamma\_policy}${=}5$ (finite, so escape is possible but not trivial).

\subsection{Outcomes and analysis}
Primary outcome: terminal mainstream conversion \param{occ\_B} (excluding the
pinned vanguard) as a function of $f$, per (topology, placement) cell. Fit each
curve for a critical fraction $f^{*}$ (the inflection of a sigmoid) and a sharpness;
$f^{*}{=}\,$undefined/linear $\Rightarrow$ simple contagion. Secondary: conversion
time; the order parameter $\mt$ trajectory; in (S), the time-course of
\param{rho\_k} on the disagreement channels (does the seal physically re-open?).

\begin{hyp}[Simple contagion]\label{h:simple}
Mainstream conversion is monotone and approximately linear in $f$ with no
threshold; topology and placement have only second-order effects.
\end{hyp}
\begin{hyp}[Complex contagion]\label{h:complex}
There is a sharp critical fraction $f^{*}$; $f^{*}$ \emph{depends on topology and
placement}---lower for clustered-behind-wide-bridge than for random or
high-centrality placement---and rises with the depth of the self-seal
(\param{cost\_scale}, $g$) and falls with $\beta_s$ and \param{q\_reliability}.
\end{hyp}

\begin{prediction}
We expect complex contagion (Hypothesis~\ref{h:complex}): a threshold near
$f^{*}\!\sim\!0.2$--$0.25$ for random placement, materially lower for clustered
wide-bridge placement, and a non-monotone (over-shooting, hence discounted)
effect for an \emph{isolated} high-centrality vanguard, consistent with the
notebook-10 isolated-vanguard overconfidence result. Higher $\beta_s$ and
\param{q\_reliability} shift $f^{*}$ down.
\end{prediction}

\severe{The deafened-vanguard control: set the mainstream's trust toward
vanguard agents to zero (zero the corresponding columns of $T$, renormalise).
Social rescue must then vanish \emph{at every $f$}. If the mainstream still
converts, the ``rescue'' is leaking through an unmodelled channel and the
trust-bandwidth-limited claim is refuted. A second falsifier: if conversion is
linear in $f$ with no topology dependence, paradigm change in this model is a
\emph{simple} contagion and the complex-contagion framing is wrong.}

% =====================================================================
\section{Shared infrastructure and reporting standards}
\label{sec:infra}

\paragraph{Order parameters (existing code).}
(S): \code{Network.run\_trace}$\to(\mt,\text{grav\_attn}_t)$;
\code{run\_trace\_bmr}$\to(\mt,\Delta F_t)$;
\code{run\_trace\_precision}$\to(\mt,\text{gw}_t,\rho_t)$.
(P): \code{run\_fast}$\to\{$\param{mean\_qB}, \param{occ\_B}, \param{p\_discrim},
\param{chosen\_d}, \param{final\_q}$\}$. No new instrumentation is required;
Experiments~1--3 are sweeps over these existing returns.

\paragraph{Statistics.}
$n_{\text{seeds}}\ge 40$ per cell; report median and inter-quartile range of every
order parameter, not just the mean (lock-in is heavy-tailed across seeds). Phase
boundaries are reported with bootstrap confidence bands over seeds. All
hypothesis tests are pre-registered one-sided tests against the stated null.

\paragraph{Multiple realisability as a corroboration criterion.}
Friston's discipline: every mechanism claim must appear in \emph{both}
implementations or be flagged as an artifact. Self-censorship (Exp.~2) must show
as a \param{chosen\_d} collapse in (P) \emph{and} a \param{rho\_k} silencing in
(S); social rescue (Exp.~3) must show as \param{occ\_B} conversion in (P)
\emph{and} a re-opening of \param{rho\_k} in (S).

\paragraph{No silent caps.}
Any truncation (max horizon, seed count, swept range) that bounds a reported
boundary must be logged next to the result, so a frontier that is really an
artifact of a too-short horizon cannot be mistaken for a structural limit.

% =====================================================================
\section{Summary: philosophy $\to$ knob $\to$ prediction $\to$ falsifier}

\begin{center}
\renewcommand{\arraystretch}{1.35}
\begin{tabular}{@{}p{2.3cm} p{3.0cm} p{3.4cm} p{3.6cm}@{}}
\toprule
\textbf{Construct} & \textbf{Knob(s)} & \textbf{Prediction} & \textbf{Falsifier} \\
\midrule
Kuhnian irreversibility vs.\ relaxation (Exp.~1)
& \param{core\_governance}, \param{gamma\_policy}, ramp rate $r$
& $A_{\text{hys}}\!\propto\! r$; zero residual loop at finite $g$; fold only at $\gamma{=}\infty$
& Residual loop $A_{\text{hys}}(r{\to}0)>0$ at any finite $g$ \\
Progressive vs.\ degenerating programme (Exp.~2)
& \param{cost\_scale}, $D_0$, \param{gamma\_policy}, \param{experiment\_bias}
& \param{chosen\_d} collapses above conviction $\sim0.7$ even at zero cost (greedy)
& At cost${=}0$, \param{chosen\_d} never collapses for any $D_0$ \\
Simple vs.\ complex contagion / committed minority (Exp.~3)
& $f$, topology \param{kind}, placement, $\beta_s$, \param{q\_reliability}
& Sharp $f^{*}\!\sim\!0.2$--$0.25$, topology-dependent; isolated central vanguard over-shoots
& Linear, topology-free conversion; or rescue survives deafened-vanguard control \\
\bottomrule
\end{tabular}
\end{center}

\paragraph{Epistemic status.}
This is a Stage-1/2 design document (forest-walk synthesis): the experiments and
their falsifiers are specified, but the predicted magnitudes are extrapolations
from notebooks 18--20 and from the cited literature, not yet established on this
exact protocol. The next step is to implement the three sweeps as notebooks 21--23
and run the falsifiers \emph{first}.

% =====================================================================
\begin{thebibliography}{9}
\bibitem{kuhn1962Structure} T.~S.~Kuhn. \emph{The Structure of Scientific
Revolutions}. University of Chicago Press, 1962. \emph{(refs.bib:
\code{kuhn1962Structure})}
\bibitem{lakatos1970methodology} I.~Lakatos. Falsification and the methodology of
scientific research programmes. In \emph{Criticism and the Growth of Knowledge},
Cambridge University Press, 1970.
\bibitem{mayo2018severe} D.~G.~Mayo. \emph{Statistical Inference as Severe
Testing: How to Get Beyond the Statistics Wars}. Cambridge University Press, 2018.
\bibitem{centola2018how} D.~Centola. \emph{How Behavior Spreads: The Science of
Complex Contagions}. Princeton University Press, 2018.
\bibitem{centola2010spread} D.~Centola. The spread of behavior in an online social
network experiment. \emph{Science}, 329(5996):1194--1197, 2010.
\bibitem{wattsStrogatz1998SmallWorld} D.~J.~Watts and S.~H.~Strogatz. Collective
dynamics of `small-world' networks. \emph{Nature}, 393:440--442, 1998.
\emph{(refs.bib: \code{wattsStrogatz1998SmallWorld})}
\bibitem{millidge2021EFE} B.~Millidge, A.~Tschantz, C.~L.~Buckley. Whence the
expected free energy? \emph{Neural Computation}, 33(2):447--482, 2021.
\emph{(refs.bib: \code{millidge2021EFE})}
\bibitem{azoulay2019StarsShadow} P.~Azoulay, C.~Fons-Rosen, J.~S.~Graff~Zivin.
Does science advance one funeral at a time? \emph{American Economic Review},
109(8):2889--2920, 2019. \emph{(refs.bib: \code{azoulay2019StarsShadow})}
\bibitem{planck1949Autobiography} M.~Planck. \emph{Scientific Autobiography and
Other Papers}. Williams \& Norgate, 1949. \emph{(refs.bib:
\code{planck1949Autobiography})}
\end{thebibliography}

\end{document}
